% GENERATED FILE. Edit canonical lessons, then run npm run generate:books.
% canonical-source-hash: fnv1a64:abe63566e94556b7
% canonical-lessons: ML-C341-vicarikkuka, ML-C341-ennu-karutunnu, ML-C341-aikyam, ML-W341-ai, ML-C341-atukondu, ML-R341-opinion-and-result

\chapter{What I Think, and So}
\label{ch:ml-what-i-think-and-so}
\hlchaptermodality{341}

\begin{chapteropening}
\textbf{By the end of this chapter:} \emph{I can say what I think with ennŭ ñān vicārikkunnu, karutunnu or enikkŭ tōnnunnu, and draw a result with atukoṇṭŭ.}

Complete ML-R341-opinion-and-result: an opinion three ways, a reason before its result with atukoṇṭŭ, and the new letter \ml{ഐ}.
\end{chapteropening}

\section[\texorpdfstring{\ml{വിചാരിക്കുക} (vicārikkuka)}{വിചാരിക്കുക (vicārikkuka)}]{\ml{വിചാരിക്കുക} (vicārikkuka) — to think, to suppose}

\label{lesson:ML-C341-vicarikkuka}



\begin{quote}
Before the new one: say the Malayalam for \textquotedblleft{}a thought\textquotedblright{}. (\emph{vicāraṁ}.) Then \textquotedblleft{}I think it is good\textquotedblright{}. (\emph{nallatāṇŭ ennŭ ñān cintikkunnu}.)
\end{quote}

\subsection*{You'll want to know: \ml{വിചാരിക്കുക}}
\textbf{\ml{വിചാരിക്കുക}} — \emph{vicārikkuka} — \textquotedblleft{}to think, to suppose\textquotedblright{}.

It is the everyday verb for what you think is so, and it takes what you think in front of it, closed with \textbf{\ml{എന്ന്}}, the way \textbf{\ml{ചിന്തിക്കുക}} did when you met it.

\begin{quote}
\textbf{\ml{ഇത്} \ml{ഒരു} \ml{ചെറിയ} \ml{പ്രശ്നമാണ്} \ml{എന്ന്} \ml{ഞാൻ} \ml{വിചാരിച്ചു}.}
\end{quote}

\begin{quote}
\emph{itŭ oru ceṟiya praśnamāṇŭ ennŭ ñān vicāriccu.} — \textquotedblleft{}I thought this was a small problem.\textquotedblright{}
\end{quote}

\begin{quote}
\textbf{\ml{അവൻ} \ml{വരും} \ml{എന്ന്} \ml{ഞാൻ} \ml{വിചാരിച്ചു}.}
\end{quote}

\begin{quote}
\emph{avan varuṁ ennŭ ñān vicāriccu.} — \textquotedblleft{}I thought he would come.\textquotedblright{}
\end{quote}

The past, \textbf{\ml{വിചാരിച്ചു}}, often carries a hint that you were wrong: he did not come after all. \textbf{\ml{ചിന്തിക്കുക}} leans the other way, towards thinking something over.

\begin{cousinweb}
\textbf{\ml{വിചാരിക്കുക}} (\emph{vicārikkuka}) is \textbf{\ml{വിചാരം}}, \textquotedblleft{}a thought\textquotedblright{}, from Sanskrit \emph{vicāra}, \textquotedblleft{}pondering\textquotedblright{} (CDIAL 11683), made into a verb with \textbf{-\ml{ഇക്കുക}}. That ending turns many Sanskrit nouns into Malayalam verbs: \textbf{\ml{സഹായം}}, \textquotedblleft{}help\textquotedblright{}, gives \textbf{\ml{സഹായിക്കുക}}, \textquotedblleft{}to help\textquotedblright{}.
\end{cousinweb}

\subsection*{Your turn}
Say these aloud:

\begin{itemize}
\raggedright
  \item \emph{vicārikkuka}, then the past \emph{vicāriccu}
  \item \textquotedblleft{}I thought this was a small problem\textquotedblright{}
  \item \textquotedblleft{}I thought he would come\textquotedblright{} — \emph{avan varuṁ ennŭ ñān vicāriccu}
  \item the noun this verb is built on — \emph{vicāraṁ}
\end{itemize}

\begin{tcolorbox}[breakable,colback=teal!4,colframe=teal!35!black,title={Before you move on}]
What does \textbf{\ml{വിചാരിക്കുക}} mean? (\textbf{to think, to suppose.}) What closes the thing you think? (\textbf{\ml{എന്ന്}}.) What hint does \emph{vicāriccu} often carry? (\textbf{That you were wrong.})
\end{tcolorbox}
\section[\texorpdfstring{\ml{എന്ന്} \ml{ഞാൻ} \ml{കരുതുന്നു} (ennŭ ñān karutunnu)}{എന്ന് ഞാൻ കരുതുന്നു (ennŭ ñān karutunnu)}]{\ml{എന്ന്} \ml{ഞാൻ} \ml{കരുതുന്നു} (ennŭ ñān karutunnu) — I think that …}

\label{lesson:ML-C341-ennu-karutunnu}



\begin{quote}
Say \textquotedblleft{}to think, to suppose\textquotedblright{}. (\emph{vicārikkuka}.) Then \textquotedblleft{}to feel, to seem\textquotedblright{}. (\emph{tōnnuka}.) Then \textquotedblleft{}an opinion\textquotedblright{}. (\emph{abhiprāyaṁ}.)
\end{quote}

\begin{grammarlens}[title={Grammar Lens: … \ml{എന്ന്} \ml{ഞാൻ} \ml{വിചാരിക്കുന്നു} / \ml{കരുതുന്നു}}]
An opinion has the same frame as a report: what you think, then \textbf{\ml{എന്ന്}}, then you, then the verb. Two verbs fill the last slot:

\begin{quote}
\textbf{\ml{ഇത്} \ml{നല്ലതാണ്} \ml{എന്ന്} \ml{ഞാൻ} \ml{വിചാരിക്കുന്നു}.}
\end{quote}

\begin{quote}
\emph{itŭ nallatāṇŭ ennŭ ñān vicārikkunnu.} — \textquotedblleft{}I think this is good.\textquotedblright{}
\end{quote}

\begin{quote}
\textbf{\ml{അത്} \ml{ശരിയാണ്} \ml{എന്ന്} \ml{ഞാൻ} \ml{കരുതുന്നു}.}
\end{quote}

\begin{quote}
\emph{atŭ śariyāṇŭ ennŭ ñān karutunnu.} — \textquotedblleft{}I consider that right.\textquotedblright{}
\end{quote}

\textbf{\ml{കരുതുക}}, \textquotedblleft{}to consider, to hold\textquotedblright{}, is a little weightier than \textbf{\ml{വിചാരിക്കുക}}: the view is thought through, not a guess.
\end{grammarlens}

\begin{grammarlens}[title={Grammar Lens: \ml{എനിക്ക്} … \ml{എന്ന്} \ml{തോന്നുന്നു}}]
The softest way puts you in the dative, as with \textbf{\ml{എനിക്ക്} \ml{അറിയാം}}: the thought comes to you.

\begin{quote}
\textbf{\ml{ഇത്} \ml{നല്ലതാണ്} \ml{എന്ന്} \ml{എനിക്ക്} \ml{തോന്നുന്നു}.}
\end{quote}

\begin{quote}
\emph{itŭ nallatāṇŭ ennŭ enikkŭ tōnnunnu.} — \textquotedblleft{}It seems to me this is good.\textquotedblright{}
\end{quote}

For a view you stand behind, open with \textbf{\ml{എന്റെ} \ml{അഭിപ്രായത്തിൽ}} (\emph{enṟe abhiprāyattil}), \textquotedblleft{}in my opinion\textquotedblright{}:

\begin{quote}
\textbf{\ml{എന്റെ} \ml{അഭിപ്രായത്തിൽ} \ml{ഇത്} \ml{ഒരു} \ml{നല്ല} \ml{തീരുമാനമാണ്}.}
\end{quote}

\begin{quote}
\emph{enṟe abhiprāyattil itŭ oru nalla tīrumānamāṇŭ.} — \textquotedblleft{}In my opinion, this is a good decision.\textquotedblright{}
\end{quote}
\end{grammarlens}

\subsection*{Your turn}
Say these aloud:

\begin{itemize}
\raggedright
  \item \textquotedblleft{}I think this is good\textquotedblright{}, with \emph{vicārikkunnu}
  \item \textquotedblleft{}I consider that right\textquotedblright{}, with \emph{karutunnu}
  \item \textquotedblleft{}It seems to me this is good\textquotedblright{}, with \emph{enikkŭ … tōnnunnu}
  \item \textquotedblleft{}In my opinion, this is a good decision.\textquotedblright{}
\end{itemize}

\begin{tcolorbox}[breakable,colback=teal!4,colframe=teal!35!black,title={Before you move on}]
Which two verbs fill \textquotedblleft{}… \emph{ennŭ ñān} \_\_\_\textquotedblright{}? (\emph{\textbf{vicārikkunnu}}, \emph{\textbf{karutunnu}}.) Which way puts you in the dative? (\emph{\textbf{enikkŭ … ennŭ tōnnunnu.}}) \textquotedblleft{}In my opinion\textquotedblright{}? (\textbf{\ml{എന്റെ} \ml{അഭിപ്രായത്തിൽ}}.)
\end{tcolorbox}
\section[\texorpdfstring{\ml{ഐക്യം} (aikyaṁ)}{ഐക്യം (aikyaṁ)}]{\ml{ഐക്യം} (aikyaṁ) — unity}

\label{lesson:ML-C341-aikyam}



\begin{quote}
Before the new one: say \textquotedblleft{}I think this is good\textquotedblright{}, then \textquotedblleft{}it seems to me this is good\textquotedblright{}.
\end{quote}

\subsection*{You'll want to know: \ml{ഐക്യം}}
\textbf{\ml{ഐക്യം}} — \emph{aikyaṁ} — \textquotedblleft{}unity, togetherness\textquotedblright{}: people holding together as one.

\begin{quote}
\textbf{\ml{ഞങ്ങളുടെ} \ml{ഗ്രാമത്തിൽ} \ml{നല്ല} \ml{ഐക്യം} \ml{ഉണ്ട്}.}
\end{quote}

\begin{quote}
\emph{ñaṅṅaḷuṭe grāmattil nalla aikyaṁ uṇṭŭ.} — \textquotedblleft{}There is real unity in our village.\textquotedblright{}
\end{quote}

\begin{quote}
\textbf{\ml{കുടുംബത്തിൽ} \ml{ഐക്യം} \ml{വേണം} \ml{എന്ന്} \ml{ഞാൻ} \ml{കരുതുന്നു}.}
\end{quote}

\begin{quote}
\emph{kuṭumbattil aikyaṁ vēṇaṁ ennŭ ñān karutunnu.} — \textquotedblleft{}I think a family needs unity.\textquotedblright{}
\end{quote}

It is a word for speeches and newspapers as much as for kitchens, and it opens with a letter you have not written yet. The next lesson puts it in your hand.

\begin{cousinweb}
\textbf{\ml{ഐക്യം}} (\emph{aikyaṁ}) $\leftarrow$ Sanskrit \emph{aikya}, \textquotedblleft{}oneness\textquotedblright{}, taken over whole as a tatsama with Malayalam's noun ending (CDIAL 2533). \emph{aikya} is built on \emph{eka}, \textquotedblleft{}one\textquotedblright{}: Sanskrit strengthens the first vowel, \emph{e} to \emph{ai}, to make \textquotedblleft{}the state of being one\textquotedblright{}. Hindi \emph{ek}, \textquotedblleft{}one\textquotedblright{}, is the same root.
\end{cousinweb}

\subsection*{Your turn}
Say these aloud:

\begin{itemize}
\raggedright
  \item \emph{aikyaṁ}, slowly, then at speed
  \item \textquotedblleft{}There is real unity in our village.\textquotedblright{}
  \item \textquotedblleft{}I think a family needs unity.\textquotedblright{}
\end{itemize}

\begin{tcolorbox}[breakable,colback=teal!4,colframe=teal!35!black,title={Before you move on}]
What does \textbf{\ml{ഐക്യം}} mean? (\textbf{unity, togetherness.}) Which Sanskrit number is inside it? (\emph{\textbf{eka}}, \textquotedblleft{}one\textquotedblright{}.)
\end{tcolorbox}
\section[\texorpdfstring{\ml{ഐ} (ai)}{ഐ (ai)}]{\ml{ഐ} (ai) — the letter that opens \ml{ഐക്യം}}

\label{lesson:ML-W341-ai}



\begin{quote}
Before the new one: write \textbf{\ml{എ}} and \textbf{\ml{ൈ}} once each, and name them. (\emph{e} at the start of a word; the \emph{ai} sign, printed before its consonant, as in \textbf{\ml{കൈ}}.)

Now say \textbf{\ml{ഐക്യം}} (\emph{aikyaṁ}), \textquotedblleft{}unity\textquotedblright{}. This lesson takes its first letter out and puts it in your hand.
\end{quote}

\subsection*{Script you'll notice: \ml{ഐ}}
\textbf{\ml{ഐ}} — \emph{ai}, at the start of a word.

After a consonant, \emph{ai} is the sign \textbf{\ml{ൈ}}, printed before the consonant it follows in speech: \textbf{\ml{കൈ}}, \emph{kai}. At the start of a word it stands as a letter of its own, \textbf{\ml{ഐ}}.

Its left half has the shape of \textbf{\ml{എ}}. The right half is a second piece: a small loop, a short base, and a tall loop that drops below the line.

\subsection*{Writing: \ml{ഐ}}
One recorded hand writes \textbf{\ml{ഐ}} in two strokes. First the left half, as you write \textbf{\ml{എ}}: from the inner tip, round the small loop and out over the arch to its foot. Then, after one lift, the right half in one run: up and over its small loop, along the base, up the short stem, down below the line, and up the tall loop, over its top and down the right side.

Write \textbf{\ml{ഐ}} yourself, slowly, and larger than it is printed. Copy the printed shape and keep the order above.

\subsection*{Your turn}
\begin{itemize}
\raggedright
  \item \emph{Find:} \ml{ഐ} at the start of \textbf{\ml{ഐക്യം}}
  \item \emph{Trace it:} \ml{ഐ} three times, saying \emph{ai} as you finish each one
  \item \emph{Write it:} \ml{എ}, then \ml{ഐ}, from memory
  \item \emph{Write it:} \textbf{\ml{ഐക്യം}}
\end{itemize}

\begin{tcolorbox}[breakable,colback=teal!4,colframe=teal!35!black,title={Before you move on}]
Which letter is this — \ml{ഐ}? (\textbf{ai}, at the start of a word.) Which word did it come from? (\textbf{\ml{ഐക്യം}}, \emph{aikyaṁ}, \textquotedblleft{}unity\textquotedblright{}.) Which letter is its left half shaped like? (\textbf{\ml{എ}}.)
\end{tcolorbox}
\section[\texorpdfstring{\ml{അതുകൊണ്ട്} \ml{ഞാൻ} \ml{വന്നില്ല} (atukoṇṭŭ ñān …)}{അതുകൊണ്ട് ഞാൻ വന്നില്ല (atukoṇṭŭ ñān …)}]{\ml{അതുകൊണ്ട്} \ml{ഞാൻ} \ml{വന്നില്ല} (atukoṇṭŭ ñān vannilla) — so I did not come}

\label{lesson:ML-C341-atukondu}



\begin{quote}
Say \textquotedblleft{}so, therefore\textquotedblright{}. (\emph{atukoṇṭŭ}.) Then \textquotedblleft{}a reason; because\textquotedblright{}. (\emph{kāraṇaṁ}.)
\end{quote}

\begin{grammarlens}[title={Grammar Lens: the reason first, then \ml{അതുകൊണ്ട്}}]
\textbf{\ml{അതുകൊണ്ട്}} is \textbf{\ml{അത്}}, \textquotedblleft{}that\textquotedblright{}, with \textbf{\ml{കൊണ്ട്}}, \textquotedblleft{}with\textquotedblright{}: \textquotedblleft{}with that\textquotedblright{}, \textquotedblleft{}so\textquotedblright{}. It points forward, from a reason to what it caused.

\begin{quote}
\textbf{\ml{എനിക്ക്} \ml{പനിയായിരുന്നു}. \ml{അതുകൊണ്ട്} \ml{ഞാൻ} \ml{വന്നില്ല}.}
\end{quote}

\begin{quote}
\emph{enikkŭ paniyāyirunnu. atukoṇṭŭ ñān vannilla.} — \textquotedblleft{}I had a fever. So I did not come.\textquotedblright{}
\end{quote}

\textbf{\ml{കാരണം}} points the other way, back from the result to its reason:

\begin{quote}
\textbf{\ml{ഞാൻ} \ml{വന്നില്ല}, \ml{കാരണം} \ml{എനിക്ക്} \ml{പനിയായിരുന്നു}.}
\end{quote}

\begin{quote}
\emph{ñān vannilla, kāraṇaṁ enikkŭ paniyāyirunnu.} — \textquotedblleft{}I did not come, because I had a fever.\textquotedblright{}
\end{quote}

Same facts, opposite order. Pick the one that puts first what your listener needs first.
\end{grammarlens}

\begin{grammarlens}[title={Grammar Lens: \ml{അതുകൊണ്ടാണ്} — that is why}]
Add \textbf{-\ml{ആണ്}} and the result is singled out: \textquotedblleft{}that is why\textquotedblright{}. The verb after it takes \textbf{-\ml{അത്}}, \textquotedblleft{}the …-ing\textquotedblright{}, so the whole sentence says \textquotedblleft{}it is with that that the …-ing happened\textquotedblright{}:

\begin{quote}
\textbf{\ml{മഴയായിരുന്നു}. \ml{അതുകൊണ്ടാണ്} \ml{ഞാൻ} \ml{വൈകിയത്}.}
\end{quote}

\begin{quote}
\emph{maḻayāyirunnu. atukoṇṭāṇŭ ñān vaikiyatŭ.} — \textquotedblleft{}It was raining. That is why I was late.\textquotedblright{}
\end{quote}

\textbf{\ml{വൈകി}} gives \textbf{\ml{വൈകിയത്}}, and \textbf{\ml{വന്നു}} gives \textbf{\ml{വന്നത്}}, the same way.
\end{grammarlens}

\subsection*{Your turn}
Say these aloud:

\begin{itemize}
\raggedright
  \item \textquotedblleft{}I had a fever. So I did not come.\textquotedblright{}
  \item the same, the other way round, with \emph{kāraṇaṁ}
  \item \textquotedblleft{}It was raining. That is why I was late.\textquotedblright{}
  \item \textquotedblleft{}I think the village needs unity. So I am going to the meeting.\textquotedblright{} — \emph{grāmattinŭ aikyaṁ vēṇaṁ ennŭ ñān karutunnu. atukoṇṭŭ ñān yōgattinŭ pōkunnu.}
\end{itemize}

\begin{tcolorbox}[breakable,colback=teal!4,colframe=teal!35!black,title={Before you move on}]
Which word points forward to a result? (\textbf{\ml{അതുകൊണ്ട്}}.) Which points back to a reason? (\textbf{\ml{കാരണം}}.) \textquotedblleft{}That is why I was late\textquotedblright{}? (\emph{\textbf{atukoṇṭāṇŭ ñān vaikiyatŭ.}})
\end{tcolorbox}
\section[(review)]{Review — what I think, and so}

\label{lesson:ML-R341-opinion-and-result}



\begin{quote}
Say \textquotedblleft{}to think, to suppose\textquotedblright{} and \textquotedblleft{}unity\textquotedblright{}. (\emph{vicārikkuka}, \emph{aikyaṁ}.) Write \textbf{\ml{ഐ}} once from memory.
\end{quote}

\subsection*{Your turn}
\begin{itemize}
\raggedright
  \item \emph{Say it:} \textquotedblleft{}I thought he would come.\textquotedblright{}
  \item \emph{Say it:} \textquotedblleft{}I think this is good\textquotedblright{}, first with \emph{vicārikkunnu}, then with \emph{karutunnu}
  \item \emph{Say it:} \textquotedblleft{}It seems to me this is good.\textquotedblright{}
  \item \emph{Say it:} \textquotedblleft{}In my opinion, this is a good decision.\textquotedblright{}
  \item \emph{Say it:} \textquotedblleft{}There is real unity in our village.\textquotedblright{}
  \item \emph{Say it:} \textquotedblleft{}I had a fever. So I did not come.\textquotedblright{} Then: \textquotedblleft{}That is why I was late.\textquotedblright{}
  \item \emph{Write it:} \textbf{\ml{ഐക്യം}}
\end{itemize}

\begin{tcolorbox}[breakable,colback=teal!4,colframe=teal!35!black,title={Before you move on}]
\textquotedblleft{}I thought (wrongly) he would come\textquotedblright{}? (\emph{\textbf{avan varuṁ ennŭ ñān vicāriccu.}}) The softest \textquotedblleft{}I think\textquotedblright{}? (\emph{\textbf{enikkŭ … ennŭ tōnnunnu.}}) Reason first, then result: which word? (\textbf{\ml{അതുകൊണ്ട്}}.) \textquotedblleft{}Unity\textquotedblright{}? (\textbf{\ml{ഐക്യം}}, which opens with \textbf{\ml{ഐ}}.)
\end{tcolorbox}
